\section{Experimental Setup}\label{sec:setup}
Our evaluation examines reward construction on LIBERO and embodiment adaptation on the CDPR. The former compares phase-aware dense feedback with terminal-only and phase-only configurations; the latter evaluates complete-task control after the proposed learning schedule. Model interfaces, scene distributions, and success verifiers are specified for each setting.

\subsection{New-Embodiment Adaptation with SmolVLA}
The adaptation study uses a cable-driven manipulator as an unconventional embodiment. The simulation implements a five-channel CDPR in MuJoCo Warp/MJ-Lab, camera inputs, four object categories (apple, orange, potato, tomato), and plate/bowl destinations. The matched evaluation uses 512 distinct development scenes, excluding the 128-scene in-run validation panel, with four repeats per checkpoint. Each episode has a 128-decision horizon and no settling decisions or teacher/yaw assistance. Deployment uses the residual mean with VLA prior noise scale 1. Numerical failures remain in all-attempt denominators. \pending{Supply simulator version, scene-generation ranges, camera calibration, and the resolved action timing.}

The main comparison evaluates the 20M and 40M latent-likelihood bounded-residual checkpoints, with a panel-selected 39.8M checkpoint as a secondary comparison. These counters measure cumulative selected actions since latent-likelihood conversion; they exclude earlier inherited training and do not count all sampled interactions. The 20M--40M leg adds 19,987,881 selected actions, 609,792 episodes, and 299 updates over 89.14 recorded training hours. The learning rate remains $10^{-5}$ and the recorded LoRA change is zero.

Reported effect intervals use a percentile bootstrap of per-scene mean differences, keeping repeats clustered by scene. Within-scene permutation $p$-values use the recorded Holm correction for the primary strict comparisons. Subgroup results are exploratory. This is one continuing training lineage; four evaluation repeats are not four independent training seeds. Final-test results and controlled comparisons of dense versus binary skill starts, self-imitation on/off, and staged versus terminal-only full-task learning remain to be added. Likelihood and prior-contribution ablations assess implementation choices separately.
% Joint-paper experimental protocol for reward construction.
\subsection{Reward Construction on LIBERO}\label{sec:libero_setup}
We test whether task-progress feedback benefits a trained policy when complete-task outcomes provide limited contrast, using the phase-aware reward in Sec.~\ref{sec:dense_method}. These experiments evaluate reward construction with an established action interface.

We train OpenVLA-OFT in SimpleVLA-RL on the ten LIBERO-Spatial tasks, using two NVIDIA L40 GPUs and an aggregate rollout batch of 32. LoRA rank and alpha are both 16, applied to the language-model projector; gradient checkpointing and FSDP offloading reduce memory demand. These updates optimize discrete action-token likelihoods, despite the continuous-action supervised formulation associated with OFT. \pending{Confirm implementation commit, tokenizer/action decoding, adapted modules, group size, and optimization schedule.}

All aggregate reward configurations start from 61.0\% validation success. The primary dense condition uses $(w_p,w_c,C,\lambda)=(0.2,0.3,0.05,5)$; its terminal-only control disables shaping and retains $\lambda=5$. The filtered dense condition additionally uses $[L,U]=[0.1,0.9]$. With $\lambda=1$, dense shaping uses $w_p=w_c=0.05$, while phase-only shaping uses $w_p=0,w_c=0.2$; both use $C=0.05$.

We report the final recorded evaluations, balanced across the ten tasks. To examine low-success learning, we restrict training and validation to \texttt{task\_1}, starting at 10\%, with $[L,U]=[0.1,1.0]$ in both conditions. \pending{Complete evaluation counts, seeds, canonical task name, elapsed budgets, and accepted-update counts.} Percentages are descriptive until these records support uncertainty estimates and budget matching.


\subsection{Physical Deployment}
\assumed{The following protocol and result prose are prepared for a completed successful experiment; retain only details confirmed by the actual trial record.}
We evaluated the selected policy on the physical cable-driven platform using \pending{hardware configuration}, \pending{camera/state-estimation setup}, and a control frequency of \pending{Hz}. The deployment condition was \pending{direct transfer, calibration only, or hardware adaptation; describe exactly}. The hardware checkpoint was \pending{path and hash}; it must not be inferred from the simulation table.

Each trial began with an empty gripper and an object outside the destination. A successful trial required acquisition, sustained lift and transport, deliberate release, and placement in the instructed receptacle without assistance. We recorded all attempts, interventions, and aborts under \pending{trial/reset protocol and stopping rule}. The evaluation comprised \pending{objects, destinations, layouts, sessions, and trials per policy}, with \pending{baseline and ordering}. The earlier outline's 80-trial design is only a suggestion, not a completed sample size.
